%========================================================== STATUS 17 SEPTEMBER
\section*{Status at 17 September --- read this before anything else}
\addcontentsline{toc}{section}{Status at 17 September --- read this before anything else}

Revision~1 of this document was written on the evening of 16 September, off site, before anyone had
touched the drone. The team is now at LC3, the race drone has been weighed, and the physical-qualifier
specification (VADR-TS-005, issue 00.02) has been re-read line by line. Some of what follows in
Parts~I and~II has been overtaken. \textbf{Where this section disagrees with a later section, this
section is right.}

\subsection*{What changed}

\begin{table}[H]\centering\footnotesize
\begin{tabularx}{\linewidth}{@{}p{3.0cm}XX@{}}
\toprule
Item & Revision~1 said & Now \\
\midrule
Gene's model & The main line for Plan~A & \textbf{Abandoned.} It does not survive on the corrected simulator. Gene has reviewed the findings and agrees. The objective is now to train a new policy on the corrected simulator \\
Drone mass & \SI{0.6076}{kg} (the simulated value) & \textbf{\SI{1745}{g}} measured, race configuration. See the box below for why this matters less, and differently, than it looks \\
Camera & fx\,$\approx$\,424\,px (\SI{74.1}{\degree}), from another team's calibration & \textbf{fx\,=\,fy\,=\,417.0\,px} (\SI{75}{\degree}), from our own specification. Ours must still be calibrated --- the spec says no calibration matrices will be provided \\
Deployment layer & ``Also still to be built'' & Still the critical path, but \textbf{much smaller than feared}: the organizers ship a complete, working MSP layer on the drone (\code{~/target/msp/}), including a fake flight controller for offline development. Ours is the adapter, not the protocol \\
Link speed & Assumed our 60\,Hz policy rate was achievable & The organizers state twice that the link gives ``attitude at 30--50\,Hz; a hard real-time control loop is not'' realistic. \textbf{May force a retrain at a lower policy rate} \\
Camera mode & Assumed 1920$\times$1080 at 60\,fps & The carrier wires only \textbf{two of four} camera data lanes. That mode may not exist --- and the \SI{75}{\degree} figure is quoted \emph{for} it \\
``Next gate'' input & Not called out & \textbf{No source exists on the real drone.} The simulator supplied it \\
Safety pilot & Assumed possibly available during a run & \textbf{A valid run allows no human pilot intervention at all} (spec \S3.3). The override switch aborts a run; it cannot rescue one \\
Track footprint & $25.9\times\SI{50.25}{m}$ from the drawing's ruler & \textbf{Resolved.} The official gate-coordinate reference gives $85\times165$\,ft $=25.91\times\SI{50.29}{m}$, confirming the drawing. The spec's ``$60\times\SI{21}{m}$'' is wrong. Our corrected track matches the official table exactly \\
Compute machine & Gene's machine & \textbf{This laptop is now the primary simulation and test machine}, with cloud GPU rented for training \\
\bottomrule
\end{tabularx}
\end{table}

\subsection*{What is still true}

The durability study (Section~\ref{sec:stress}) stands, and its ranking is the most useful thing in
this document: \textbf{command delay and thrust calibration dominate every other error source}, with
camera field of view and camera tilt next. The gate-counting bug analysis (Section~\ref{sec:bug})
stands and the fix is applied. The map-audit method stands. Plan~B's logic and speed budget stand,
with the caveat that its numbers come from a purpose-built simulator, not from Isaac.

\begin{warnbox}[The weight finding, stated carefully --- this one is easy to get wrong]
The real drone is \SI{1745}{g}; the simulated one is \SI{607.6}{g}, an 8-inch airframe against a
5-inch one. The obvious response --- ``change the mass in the simulator and retrain'' --- is
\textbf{almost entirely wasted effort}, and doing it carelessly makes things worse.

The simulator converts a throttle command into a force by multiplying by the drone's own mass:
\code{thrust = (stick / hover\_stick) $\times$ mass $\times$ g}. Acceleration is force divided by
mass, so the mass cancels out. Change only the mass and the simulated drone flies identically.

Three things do \emph{not} cancel, and these are the real issues:
\begin{enumerate}
\item \textbf{Thrust-to-weight.} The simulator assumes the drone hovers at 25.5\,\% throttle, which
implies it can produce 3.5 times its own weight in thrust. That is a \SI{600}{g} racing quad's
figure. A \SI{1745}{g} machine carrying a Jetson more likely hovers near 35--45\,\%, i.e.\ about
2.0--2.6 times its weight. If so, the simulator has been promising the policy acceleration the drone
does not have. \textbf{Measuring this takes five minutes} (hover, read the throttle) and it is the
single most valuable number available today. For scale: a \SI{10}{\%} thrust error on its own took
crashes from about 2 to \textbf{58 per 100 gates} in the durability study.
\item \textbf{Rotational inertia.} The real airframe resists being rotated roughly \textbf{11--16
times} more than the simulated one (2.9$\times$ the mass, with arms about 2.2$\times$ longer, and
inertia grows with the square of the arm). \textbf{Danger:} the simulator's rate loop uses a fixed
gain \code{rate\_kp = 0.08}. If someone sets the real inertia and leaves that gain alone, the
simulated drone becomes about 13$\times$ more sluggish than anything real, and every policy trained
on it is worthless. That gain is not physics --- it stands in for Betaflight's own controller, which
is retuned for whatever airframe it is on. It must be re-derived from a measured response:
\code{gain $\approx$ inertia / measured response time}.
\item \textbf{Air drag}, which the simulator does not model at all.
\end{enumerate}
\end{warnbox}

\subsection*{The two gaps that block everything}

\begin{enumerate}
\item \textbf{Our flight client cannot talk to the real drone.} Searching the entire flight client
for ``MSP'' returns nothing; \code{controller.py} speaks MAVLink to the virtual qualifier's
simulator, which the real flight controller does not understand. \textbf{Much mitigated on 17
September:} the organizers ship a complete MSP layer on the board (\S\ref{sec:provided}), so what
remains is an adapter, the Betaflight curve inversion and the sign/axis mapping --- hours rather than
days, and testable on a laptop against their fake flight controller. Still the critical path.

\item \textbf{Nothing tells the drone which gate is next.} The policy's input includes a list of 18
on/off flags meaning ``this is the gate you are going for''. In the simulator that came free, from
the simulator. On the real course we must count gates ourselves from the camera --- and the double
gate counts twice per lap, which is exactly where a counter loses its place.
\end{enumerate}

\begin{warnbox}[New on 17 September: the link may be slower than our policy]
The organizers' software guide states, twice and emphatically, that the flight-controller link is
\emph{telemetry-grade}: ``MSP is request/response over a 115200-baud serial line. You poll; the FC
answers. Attitude at 30--50\,Hz is realistic; a hard real-time control loop is not.''

Our policy runs at \textbf{60\,Hz} and takes gyro readings as part of its input. Sending stick
commands at 60\,Hz is fine --- that is exactly what a radio link does, and Betaflight keeps the fast
stabilisation loop internally, which is where it belongs. The problem is the \emph{feedback} half: if
gyro only arrives at 30--50\,Hz, the simulator has been training the policy against a sensor faster
than the one it will fly with.

\textbf{Action:} measure the achievable rate on day one with the organizers' own
\code{companion\_listener\_msp.py}, then set the simulator's decimation to match \emph{before}
committing to a long training run. At 120\,Hz physics, decimation 2 gives 60\,Hz, 3 gives 40\,Hz and
4 gives 30\,Hz.
\end{warnbox}

\subsection*{What the organizers already provide}\label{sec:provided}

Everything below is pre-installed in \code{~/target/} on the Jetson. No pip, no network, no root.
Log in over a USB cable: \code{ssh dcl@192.168.55.1}, password \code{dcl}.

\begin{itemize}\sloppy
\item \code{msp/msp.py} --- \code{MSPLink}: all MSP framing, plus \code{attitude()},
\code{raw\_imu()}, \code{analog()}, \code{rc\_channels()}, \code{motors()}, \code{status()} and
arming-blocker decoding. Thread-safe, background receive loop, counts corrupt frames rather than
crashing on them.
\item \code{msp/msp\_rc.py} --- \code{RCTransmitter}: streams stick commands at a fixed rate from a
background thread \textbf{with a watchdog already built in}. Stop calling \code{set\_control()} and
it falls back to an idle frame by itself. \code{arm()} refuses unless throttle is at minimum;
\code{close()} disarms cleanly.
\item \code{msp/fake\_fc.py} --- \textbf{simulates a Betaflight flight controller on a laptop}, with
deliberate fault modes (frozen, saturated, wrong scale, no accelerometer, bus errors). The whole test
suite runs with nothing plugged in.
\item \code{msp/setup\_jetson\_uart.sh} --- frees the serial port from Linux's login console.
\textbf{Run once per board before anything else}; skipping it looks exactly like a dead flight
controller.
\item \code{msp\_bench.py} (info / telemetry / rc / demo), \code{imu\_check.py} (a real functional
IMU test, not just ``did it reply''), \code{companion\_listener\_msp.py} (live telemetry, and it
reports the achievable poll rate).
\item Camera: \code{live-view.py}, \code{live-view-pts.py} (hardware capture timestamps),
\code{live-view-imu.py}, \code{show-camera.py}, \code{raw-view.py}, \code{frame-timestamps.py}, plus
\code{bringup-check.sh} and \code{camera-bind-check.sh}.
\end{itemize}

The guide says it plainly: \textbf{``do not re-implement MSP framing.''}

\textbf{First action for whoever reaches a board:} \code{scp -r dcl@192.168.55.1:~/target ./}, and
share it --- that is what lets the rest of the team write and test adapter code without the drone.

\begin{warnbox}[Three constraints from those guides that will otherwise cost an afternoon each]
\begin{enumerate}
\item \textbf{Only two of the camera's four data lanes are wired} on this carrier, capping resolution
$\times$ framerate. Run \code{v4l2-ctl -d /dev/video0 -{}-list-formats-ext} and believe that list, not
the datasheet. The \SI{75}{\degree} field of view is quoted for 1920$\times$1080 at 60\,fps.
\item \textbf{The image processor needs a graphics context.} Colour, auto-exposure and white balance
run in NVIDIA's ISP, which fails over a plain SSH session and falls back to a flat, dim software
debayer. That is documented behaviour, not a broken camera.
\item \textbf{No deep-learning framework is installed.} OpenCV (GStreamer-enabled), NumPy, pyserial
and python3-gi --- no PyTorch. This makes the NumPy policy runtime \emph{mandatory}. And never
\code{pip install opencv-python}: it silently replaces the GStreamer build and the camera pipeline
then fails with no error text.
\end{enumerate}
There is also \textbf{no hardware sync wire between the camera and the flight controller}. Their
guide calls its camera--IMU timing section ``the section most likely to save you a week''. It is
right; read it before building anything that fuses the two.
\end{warnbox}

\subsection*{Where the working documents are}

This PDF is the onboarding and history document. The live operational plan is kept as Markdown in
the workspace root, because it changes daily:
\code{FIELD\_PLAN.md} (workstreams, procedures, schedule, contingencies),
\code{SIM\_MEASUREMENTS.md} (every parameter the simulator guesses and how to measure it),
\code{INSTALL\_PLAN.md} (what to install and download, ordered by how badly it blocks us),
\code{PLAN\_B.md} (the slow fallback).
Sections~\ref{sec:gaps}--\ref{sec:schedule} of this document summarise them.
